The clearest result of the study concerned allocation. On real prices without the market hedge (when on, the hedge
capped the portfolio's beta at each re-plan, at a level that was lower when the market was weaker, and let up to five stocks with zero or
negative beta enter), the seven weighting rules applied to one selection of $33$ stocks (three in each of the $11$
sectors) ended within $0.15$ of one another, and all seven ended below SPY. The selections themselves ended further
apart, from $1.03$ to $1.70$ times initial wealth across all combinations, and the wide spread came from the
$13$-stock selections, made across the whole universe of $373$ stocks without a sector quota, and from the two lists the author picked in
advance ($13$ stocks, and the same $13$ plus two). One explanation, consistent with these runs but not tested by them,
is that the shared limits left the weighting rules little room. Every rule, including the three benchmark portfolios (minimum variance, tangency and $130/30$ long-short), faced the
same cap on each stock's weight and the same cap on portfolio volatility, so the rules differed mainly in how they
spread the invested wealth among the selected stocks within those caps. The results do not show that any one
selection rule is reliably better on real data. The hand-picked $13$-stock list ended highest, but it is one list chosen by one
person, and its result cannot be separated from chance. The three score-based $13$-stock rules, which took the $13$
highest scores (the expected growth rate from a single-index fit divided by a power of the stock's market beta that
depended on the market's recent direction) with no screen, a correlation screen or a liquidity screen, ended above
the two random $13$-stock controls, drawn from the whole universe or from its $100$ most liquid stocks, under every
weighting rule, while the three $33$-stock rules (ranked by score, chosen by a
bandit, or drawn at random within each sector) ended in overlapping ranges. The $13$-stock rules were added after the
backtests of the $33$-stock rules on the same trading window had been seen, so their comparison with random
selections is exploratory, and one historical path cannot show whether the difference would
persist~\cite{baileyEtAl2014, harveyLiu2015}.

The rebalancing controller re-planned, re-solving the allocation from the current holdings, on a schedule set partly
by its own forecast errors. Across combinations of selection and weighting rule, the median of each combination's
$30$ runs brought a fifth to three fifths of its re-plans forward of the $21$-day
horizon, through a band exit (the realized portfolio value leaving the range the forecast considered likely) or a
drawdown stop (wealth falling $8\%$ below its peak since the last drawdown stop), and early re-plans were more frequent for stocks selected by
score than for random stocks, a difference the study did not examine. The forecast, a geometric Brownian motion,
represents neither the heavy tails nor the volatility clustering of daily growth rates~\cite{mandelbrot1963,
cont2001}, so band exits are expected to occur more often than they would if the forecast model were correct, and
the band served both as a trigger and as a running check on the forecast. A forecast model that represents
volatility clustering~\cite{engle1982, bollerslev1986} would widen the band after volatile days and narrow it after
calm ones, so its effect on the number of early re-plans cannot be signed in advance, and this was not tested.
Against re-planning every $21$ trading days alone, the early re-plans made the median worst drawdown smaller in most
combinations of selection and weighting rule, with the hedge on or off, and raised trading cost in all of them.
Median terminal wealth, by contrast, rose with the early re-plans in about half of the combinations with the hedge
off and in fewer than half with it on. The same pattern held in the two stress sets of synthetic futures, whose
declines were deeper than the real one: the early re-plans made the median worst drawdown smaller in nearly every
combination, by about as much as on real prices, and raised median terminal wealth in fewer than half of the
combinations in the set with jump bursts and in most of those in the set without (Tables~\ref{tab:S-calendar-burst} and~\ref{tab:S-calendar-bear}).
The triggers therefore limited losses more than they added growth, and the protection they gave did not grow with
the depth of the decline.

% results/synthetic/tables/{burst,bear}_stress_hedge_{off,on}_wealth_median.csv (lowest cells 0.92412 and 0.77386
% against equal-weight 0.91555 and 0.75313, margins 0.0086 and 0.0207); results/synthetic/rows/path_*.csv mean_portfolio_beta, ten main selection rules, hedge
% off, median over seeds within a path, then over the 200 paths of each stress set (burst 0.322--0.764, bear
% 0.229--0.850); real-price range from results/*_summary.csv as in results.tex
Results on the simulated markets describe how this system behaved in one generator's markets and are not forecasts.
In both stress sets, the futures with the deepest declines of simulated SPY, one drawn from futures whose market path
contained a jump episode (a sustained burst of large daily moves) and one from futures without, every combination of
selection and weighting rule ended above simulated SPY bought and held in the median, with the hedge on or off.
Random selections did so as well, so choosing stocks by score or by bandit was not needed for this result, and every
combination also did so with the band exit and the drawdown stop switched off, so the early re-plans were not needed
either (Tables~\ref{tab:S-calendar-burst} and~\ref{tab:S-calendar-bear}). The equal-weight universe,
which has neither limits nor triggers, also ended above simulated SPY in both stress sets, so part of the margin over
SPY belonged to the simulated market itself, and every combination's median exceeded the equal-weight median by as
little as $0.009$ in the set with jump episodes and $0.021$ in the set without. A likely cause of the margin over SPY is
that the portfolios carried less market exposure than the index. With the hedge off, the average market beta of the
portfolios was below one in every combination, at $0.32$ to $0.76$ in the stress set with jump episodes and $0.23$ to
$0.85$ in the set without (medians over futures), as on real prices ($0.39$ to $0.90$), and the market term of a
portfolio with a beta below one falls less than the market when the market factor falls, whatever stocks it holds.
The volatility cap, set at SPY's historical level, is one reason the beta stays below one, but the study did not run
the system without it.

The synthetic market also lacked two features of real sell-offs. Each stock's own simulated residual was generated
independently of the others, so co-movement among stocks came only from the simulated SPY and the sector factors.
Real residual co-movement tends to rise in sell-offs~\cite{longinSolnik2001, angChen2002}, and its absence may
understate how much a portfolio of a few stocks loses relative to the index. The generator's volatility clustering
also decayed within about a day, faster than in the data, so turbulent stretches in the synthetic futures were
shorter than in real markets, which may affect both the depth of drawdowns and how often the triggers fired. The
effect of the market hedge also changed sign across markets. The hedge lowered
market exposure in every combination on real prices but reduced terminal wealth in most of them, and it reduced
terminal wealth in most combinations of the $600$ typical futures, which were drawn at random from $10{,}000$
simulated paths. In the stress set without jump episodes it raised terminal wealth in most combinations, and in the
stress set with jump episodes it raised terminal wealth in about half of them. Since the hedge's beta cap tightens after
the market has started to fall, and takes effect only at the next re-plan, it pays only when the decline continues.
Between re-plans the shares were held, so the portfolio's beta was not controlled, and in a slow decline that stayed
inside the forecast band and above the drawdown limit the book could keep the shares of its last re-plan for up to
$21$ trading days.

The comparison of score-based and random selections on synthetic futures needs a further qualification. The
generator gave each synthetic stock a constant intercept, its 2014--2024 single-index intercept shrunk toward the
average across stocks. In the synthetic market, a stock's intercept on the calibration years was therefore positively
correlated with its intercept on the trading window: across stocks, the two correlated at $0.10$, against $-0.12$ in the real market.
The score's trend term is a daily updated estimate of this intercept, so in the synthetic futures it carried
information about future growth by construction, and the share of typical futures in which score-based selections
beat random ones is partly a consequence of that design. The spread of estimated intercepts across stocks was also
wider in the synthetic market than in the real one ($0.089$ against $0.076$ per year), and a wider spread of
persistent intercepts gives a score-based selection more to find. The intercept spread was the one check the
generator failed among the $22$ that had a pass criterion, whereas the intercept correlation and the clustering
shortfall were reported without one. On real prices, the score-based $13$-stock rules still ended above random
selections even though intercepts did not persist from 2014--2024 into 2025--2026. The study does not separate
whether that came from the score's market term, from shorter-lived trends picked up by the daily updated estimate,
or from chance.

The study has further limits. The trading window covered $420$ trading days with one broad sell-off, in the spring
of 2025. The universe held only stocks that were members of the S\&P~500 index at the date of the constituent
snapshot and traded on every day from 2014 to 2026, so it carries survivorship bias~\cite{brownEtAl1992}. Neither
the portfolios nor SPY earned dividends, so every comparison is in prices alone. The simulated fills charged half the
bid-ask spread and a linear market-impact term, the turnover limit, a cap on the half-spread cost of the solver's share changes at each re-plan, used the half spread alone, and taxes were not
modeled. A paper-trading deployment of the system is running, but its results are not part of this study. Natural
extensions are a forecast model with heavy tails and volatility clustering for the controller, runs without the volatility cap to measure its share of the stress-set results, residual
co-movement across stocks and less persistent intercepts in the synthetic market, and evaluation of selection rules
over many historical windows.
